\exercisegroup

\begin{enumerate}
\def\labelenumi{\arabic{enumi})}
\item
  Let X be a path-connected space, let x be a point of X, and let \(\mathcal{U} = (\mathrm{U}_{i})\) be an open covering of X; for every i, let \(c_{i}\) be the class of a path in X with origin a point \(u_i\) of \(U_i\) and ending at \(x\) . Denote by \(G_{\mathcal{U}}\) the smallest normal subgroup of \(\pi_1(X,x)\) such that \(\text{Int}(c_i)(G_{\mathcal{U}})\) contains the class of every loop in \(U_i\) at \(u_i\) . Then \(G_{\mathcal{U}}\) is an open subgroup of \(\pi_1(X,x)\) for the admissible topology.
\item
  Let X be a topological space and let \(a\) be a point of X. We say that a subgroup H of \(\pi_1(\mathrm{X}, a)\) is adequate if, for all \(b \in \mathrm{X}\) and every path class \(\gamma \in \varpi_{a,b}(\mathrm{X})\) , there exists a neighborhood V of \(b\) such that for every loop \(c \in \Omega_b(\mathrm{V})\) the class of the loop \(\gamma[c]\gamma^{-1}\) belongs to H.
\end{enumerate}

\begin{enumerate}
\def\labelenumi{\alph{enumi})}
\item
  Show that there exists a unique group topology on \(\pi_{1}(X,a)\) for which the adequate subgroups are the open subgroups.
\item
  Show that this topology is finer than the topology of compact convergence, but that these two topologies have the same open normal subgroups.
\item
  If X is locally path-connected, the adequate topology and the topology of compact convergence have the same open subgroups.
\end{enumerate}

\begin{enumerate}
\def\labelenumi{\arabic{enumi})}
\setcounter{enumi}{2}
\tightlist
\item
  Let C be the circle with diameter \(((0,0),(0,1))\) in \(R^{2}\) and let P be the union of the \(2^{n}C\) , for \(n \in Z\) . Let E be the quotient of the space \(P \times R\) by the coarsest equivalence relation for which \((x,t) \sim (2x,t+1)\) for all \((x,t) \in P \times R\) . Let p denote the canonical surjection from \(P \times R\) onto E. Let a be the point \((0,0)\) of P and set \(e = p(a,0)\) .
\end{enumerate}

\begin{enumerate}
\def\labelenumi{\alph{enumi})}
\item
  Show that p makes \(P \times R\) a Galois covering with group Z. Show that \(p_{*}(\pi_{1}(P \times \mathbf{R}, (a,0)))\) is an open normal subgroup of \(\pi_{1}(E, e)\) with quotient Z.
\item
  Let \(\varphi\colon\mathrm{P}\to\mathrm{P}\) be the map defined by \(\varphi(x)=x\) if \(x\in\mathrm{C}\) and \(\varphi(x)=(0,0)\) otherwise. Show that it is continuous.
\item
  Show that the kernel of the homomorphism \(\varphi_{*} = \pi_{1}(\varphi, (a,0))\) is a subgroup of \(\pi_{1}(\mathrm{P} \times \mathbf{R}, (a,0))\) which is open for the topology of compact convergence.
\item
  Let \(u: I \to C\) be the loop given by \(t \mapsto \frac{1}{2}(1 - \cos(2\pi t), \sin(2\pi t))\) and let \(c \in \Omega_e(E)\) be the loop given by \(t \mapsto (u(t), 0)\) . Show that the image of \(\varphi_*\) is an open subgroup of \(\pi_1(E, e)\) which does not contain the class of the loop c.
\item
  Let \(c_{n}: I \to E\) be the loop defined by \(c_{n}(t) = p(2^{-n}u(t), 0)\) and let \(d: I \to E\) be the loop given by \(t \mapsto p(a, t)\) . Show that \(d^{n}c_{n}\overline{d}^{n}\) is strictly homotopic to c. Deduce that the class \([c]\) is contained in every admissible subgroup of \(\pi_{1}(E, e)\) .
\item
  Show that the admissible topology on \(\pi_{1}(\mathrm{E},e)\) is strictly coarser than the topology of compact convergence.
\end{enumerate}

\begin{enumerate}
\def\labelenumi{\arabic{enumi})}
\setcounter{enumi}{3}
\tightlist
\item
  Let L be the Alexandroff half-line (TG, IV, p. 49, exerc. 12) and let \(L^{*}\) be its Alexandroff compactification; let \(\omega\) be the unique point of \(L^{*} - L\). Let A be a set with at least two elements and let a be a point of A; set \(X = L^{*} \times A\) and \(Y = X - \{(\omega, a)\}\).
\end{enumerate}

\begin{enumerate}
\def\labelenumi{\alph{enumi})}
\item
  Show that the canonical injection j from Y into X is étale, separated, and has the path-lifting property.
\item
  Show that the map j does not make Y a covering of X.
\end{enumerate}

\begin{enumerate}
\def\labelenumi{\arabic{enumi})}
\setcounter{enumi}{4}
\tightlist
\item
  Let \(X_{0}\) be the unit circle of the numerical plane \(R^{2}\) ; let \(X_{1}\) be the union of the segments joining the point \((1,0)\) to the points \((3,1/n)\) , for \(n \in N^{*}\) ; let \(X_{2}\) be the set of points \((x,y)\) of the plane such that \((x-2)^{2} + y^{2} = 1\) and \(y \leqslant 0\) ; we set \(X = X_{0} \cup X_{1} \cup X_{2}\) .
\end{enumerate}

\begin{enumerate}
\def\labelenumi{\alph{enumi})}
\item
  Let \(Y = X \cup (-X)\) . Let \(p: Y \to X\) be the map defined as follows: we have \(p(\cos(t), \sin(t)) = (\cos(2t), \sin(2t))\) for \(t \in R\) and \(p(z) = p(-z) = z\) for \(z \in X_{1} \cup X_{2}\) . Show that p makes Y a covering of degree 2 of X.
\item
  Let \(Z_{2}\) be the set of points of the plane such that \((x-1)^{2}+y^{2}=4\) and \(y\leqslant0\) ; we set \(Z=X_{0}\cup X_{1}\cup(-X_{1})\cup Z_{2}\cup(-Z_{2})\) . Let \(q:Z\to X\) be the map defined by \(q(\cos(t),\sin(t))=(\cos(2t),\sin(2t))\) for \(t\in R\) , \(q(z)=q(-z)=z\) for \(z\in X_{1}\) and \(q(z)=q(-z)=1+\frac{1}{2}z\) for \(z\in Z_{2}\) . Show that q makes Z a covering of degree 2 of X.
\item
  Prove that \(q_{*}\pi_{1}(\mathrm{Z},(1,0)) = p_{*}\pi_{1}(\mathrm{Y},(1,0))\) but that there does not exist a continuous map \(f\colon Z \to Y\) such that \(p \circ f = q\). \(^{(4)}\)
\end{enumerate}

\begin{enumerate}
\def\labelenumi{\arabic{enumi})}
\setcounter{enumi}{5}
\tightlist
\item
  For every integer \(n \geqslant 1\) we set \(x_{n} = (1/n,0)\) and denote by \(\mathrm{C}_{n}\) the circle of the numerical plane with diameter \((0,x_{n})\) . Let P be the union of the spaces \(\mathrm{C}_{n}\) , for \(n \in \mathbf{N}\) ("Hawaiian earring", cf. I, p. 146, exerc. 5); for every integer \(n \geqslant 1\) , let \(\mathrm{P}_{n}\) be the union of the spaces \(\mathrm{C}_{m}\) , for \(1 \leqslant m \leqslant n\) and \(\mathrm{P}^{n}\) the union of the spaces \(\mathrm{C}_{m}\) , for \(m > n\) .
\end{enumerate}

\begin{enumerate}
\def\labelenumi{\alph{enumi})}
\item
  Let n be an integer \(\geqslant 1\) and let \(r_{n}\) be the map from P to \(P_{n}\) which maps a point \(x \in P\) to itself if \(x \in P_{n}\) , and to 0 otherwise. Show that \(r_{n}\) is a continuous retraction of the canonical injection of \(P_{n}\) into P.
\item
  Prove that \(P_{n}\) is locally contractible. Deduce that every point of P distinct from the origin has a contractible neighborhood at that point.
\item
  Show that the canonical inclusions of \(P_{n}\) and \(P^{n}\) into P induce a group isomorphism of the free product \(\pi_{1}(\mathrm{P}_{n},0)*\pi_{1}(\mathrm{P}^{n},0)\) onto \(\pi_{1}(\mathrm{P},0)\) . (Show that the canonical injection of \(P^{n}\) into the complement \(U_{n}\) of the set \(\{x_{1},\ldots,x_{n}\}\) is a homotopy equivalence. Then apply proposition 2 of IV, p. 422 to the covering \((\mathrm{P}_{n},\mathrm{U}_{n})\) of P.)
\end{enumerate}

\(d)\) Let \(\gamma \in \pi_1(\mathrm{P},0)\) be a class of loops such that \((r_n)_*(\gamma) = \varepsilon_0\) for every integer \(n \geqslant 1\) ; prove that \(\gamma = \varepsilon_0\) .

\begin{enumerate}
\def\labelenumi{\alph{enumi})}
\setcounter{enumi}{4}
\item
  The canonical homomorphism of \(\pi_1(\mathrm{P},0)\) into \(\varprojlim_{n}\pi_{1}(\mathrm{P}_{n},0)\) deduced from the family \(((r_n)_*)\) is injective and the admissible topology on \(\pi_1(\mathrm{P},0)\) is the inverse image of the topology of the projective limit on \(\varprojlim_{n}\pi_{1}(\mathrm{P}_{n},0)\) , the groups \(\pi_1(\mathrm{P}_n,0)\) being equipped with the discrete topology.
\item
  Show that the path-connected components of \(\Omega_{0}(\mathrm{P})\) are closed, as are those of \(\Omega_{(0,0)}(\mathrm{P} \times \mathrm{P})\) .
\end{enumerate}

\begin{enumerate}
\def\labelenumi{\arabic{enumi})}
\setcounter{enumi}{6}
\tightlist
\item
  We resume the notation of exercise 6.
\end{enumerate}

\begin{enumerate}
\def\labelenumi{\alph{enumi})}
\item
  For every integer \(n \geqslant 1\) , let \(\alpha_{n} \in \pi_{1}(P,0)\) be the class of a loop in \(C_{n}\) whose class generates \(\pi_{1}(C_{n},0)\) . Let us then set \(\beta_{n} = \alpha_{1}\alpha_{n}\alpha_{1}^{-1}\alpha_{n}^{-1}\) and \(\gamma_{n} = \beta_{n}^{n}\) . The sequence \((\gamma_{n})\) tends to \(\varepsilon_{0}\) for the admissible topology.
\item
  Let \(n \in N^{*}\) . For all \(c \in \Omega_{0}(P)\) , let \(\omega_{n}(c)\) be the supremum of the integers k such that there exists a strictly increasing sequence \((t_{1}, \ldots, t_{2k})\) of elements of I such that \(c(t_{2i}) = 0\) and \(c(t_{2i-1}) = x_{n}\) for \(i \in \{1, \ldots, k\}\) . Prove that \(\omega_{n}(c)\) is an integer. Show that the map \(\omega_{n}\) of \(\Omega_{0}(P)\) into N is upper semicontinuous.
\item
  For every integer \(n \geqslant 1\) , let \(c_{n}\) be a path whose class is \(\gamma_{n}\) . Prove that \(\omega_{1}(c_{n}) \geqslant 4n\) . Deduce that the sequence \((c_{n})\) has no limit in \(\mathcal{C}_{c}(\mathbf{I};\mathbf{P})\) . d) The admissible topology on the group \(\pi_{1}(\mathrm{P},0)\) is Hausdorff, non-discrete, and strictly coarser than the topology of compact convergence. *The space P is in particular not delaceable.*
\end{enumerate}

\begin{enumerate}
\def\labelenumi{\arabic{enumi})}
\setcounter{enumi}{7}
\tightlist
\item
  We resume the notation of exercises 6 and 7; let also \(p\colon \Omega_{0}(\mathrm{P})\to\pi_1(\mathrm{P},0)\) be the canonical surjection.
\end{enumerate}

\begin{enumerate}
\def\labelenumi{\alph{enumi})}
\item
  For every pair \((m,n)\) of integers \(\geqslant 1\) such that \(m \neq n\) we set \(\lambda_{m,n} = (\alpha_{m}\alpha_{n}\alpha_{m}^{-1}\alpha_{n}^{-1})^{m+n}\) and \(\lambda_{m,n} = (\alpha_{1}\alpha_{m}\alpha_{1}^{-1}\alpha_{m}^{-1})^{n}\) ; let S be the set of pairs of the form \((\lambda_{m,n},\mu_{m,n})\) in \(\pi_{1}(P,0)\times\pi_{1}(P,0)\) . Show that S is not closed but that its inverse image in \(\Omega_{0}(P)\times\Omega_{0}(P)\) is. Deduce that the map \(p\times p\colon\Omega_{0}(P)^{2}\to\pi_{1}(P,0)^{2}\) is not strict.
\item
  Let T be the set of loop classes of the form \(\lambda_{m,n}\mu_{m,n}\) in \(\pi_1(\mathrm{P},0)\) . Show that T is closed.
\item
  Show that the composition law of \(\pi_{1}(P,0)\) is not continuous when the group \(\pi_{1}(P,0)\) is equipped with the quotient topology of compact convergence. \(^{(5)}\)
\end{enumerate}

\begin{enumerate}
\def\labelenumi{\arabic{enumi})}
\setcounter{enumi}{8}
\tightlist
\item
  Let us take up the notation of Exercise 6. For every integer \(n \in \mathbf{N}^*\) , let \(c_n: \mathbf{I} \to \mathrm{C}_n\) be the loop at the origin defined by \(t \mapsto \frac{1}{2n}(1 - \cos(2\pi t), \sin(2\pi t))\) .\\
\end{enumerate}

\begin{enumerate}
\def\labelenumi{\alph{enumi})}
\item
  For any sequence \(\alpha \in \{0,1\}^{\mathbf{N}^*}\) , let \(c_\alpha\) be the map defined by \(c_\alpha(0) = 0\) and \(c_\alpha(t) = \alpha_n c_n (2^n t - 1)\) for every integer \(n \geqslant 1\) and every \(t \in \mathbf{I}\) such that \(2^{-n} \leqslant t \leqslant 2^{1-n}\) . Show that this is a loop at the origin. Show that the map \(\alpha \mapsto [c_\alpha]\) of \(\{0,1\}^{\mathbf{N}^*}\) into \(\pi_1(\mathrm{P},0)\) is injective.
\item
  For every integer \(n \geqslant 1\) , let \(q_{n}\) be the map from P to \(C_{n}\) given by \(q_{n}(x) = x\) if \(x \in C_{n}\) and \(q_{n}(x) = 0\) otherwise; it is continuous. Show that the group homomorphism \(\psi: \pi_{1}(\mathrm{P}, 0) \to \prod_{n \geqslant 1} \pi_{1}(\mathrm{C}_{n}, 0)\) given by \(\psi(\gamma) = ((q_{n})_{*}(\gamma))_{n}\) admits a section.
\end{enumerate}

\begin{enumerate}
\def\labelenumi{\arabic{enumi})}
\setcounter{enumi}{9}
\tightlist
\item
  For every integer \(n \geqslant 1\) , let \(\mathrm{C}_n\) be the circle in the number plane with center \((1/n,0)\) and of radius \(1/n\) ; let P be the union of the spaces \(\mathrm{C}_n\) , for \(n \geqslant 1\) . For every integer \(n \geqslant 1\) , let \(\mathrm{B}_n\) be the set of points \((x,y)\) of the plane such that \((nx-1)^2 + n^2y^2 \leqslant 1 \leqslant ((n+1)x-1)^2 + (n+1)^2y^2\) ; let B be the set of points \((x,y)\) of the plane such that \((x-1)^2 + y^2 \leqslant 1\) . Set \(\mathrm{B}_0 = \mathrm{P}\) and let A be the quotient space of the sum space of the family \((\mathrm{B}_n)_{n \geqslant 0}\) by the coarsest equivalence relation identifying \((b,n)\) and \((b,0)\) if \(b \in \mathrm{C}_n \cup \mathrm{C}_{n+1}\) .
\end{enumerate}

\begin{enumerate}
\def\labelenumi{\alph{enumi})}
\item
  Prove that there exists a unique continuous map \(p \colon \mathrm{A} \to \mathrm{B}\) which, for every integer \(n \geqslant 0\) and every \(b \in \mathrm{B}_n\), maps the class of \((b, n)\) to the point \(b\) . Prove that \(p\) is bijective but is not a homeomorphism.
\item
  Prove that the space A is locally path-connected.
\item
  Let a be the class of the point \((0,0)\) . Prove that the quotient topology of the topology of compact convergence on \(\pi_{1}(A,a)\) is the coarse topology.
\item
  Prove that the space A is simply connected.
\item
  Let j be the canonical map from \(C = B_{0}\) into A. Show that the map \(j_{*} \colon \pi_{1}(C,0) \to \pi_{1}(A,j(0))\) is surjective and that its kernel has cardinality continuum.
\item
  Prove that the group \(\pi_{1}(A,j(0))\) has cardinality continuum. *In particular, the space A is not delaceable.*
\end{enumerate}

